\SourceLeaf{155}
whose study can be of keen interest only to someone who
already knows how to use the methods. I usually present three
justifications.

\begin{enumerate}
\item[1)] one based solely on principles of expediency, which can easily be given in a very
convincing way,
\item[2)] the way taught in the Th. M. C. C. of connecting the method
with probability theory,
\item[3)] the \textls{entirely different} way, presented in the Theoria Combinationis
Observationum, of connecting the method with probability theory. I am convinced
that it is the sole
admissible one.
\end{enumerate}

For anyone still entirely unfamiliar with this subject, I consider this
order the most expedient. I know of no expositions by others, however,
that I could recommend. As for No. 1), Mr
\textsc{Warensdorf} can learn it from countless books (as far as the A. B. C. of the matter
is concerned), and it makes little difference which he chooses, provided he later
takes the opportunity to improve for himself what almost all of them do poorly.
For example, these books do not place sufficient emphasis on the fact
that even when the observed quantities are already linear
functions of the unknown elements (as with the pendulum length
$x+y\sin\varphi^2$), one should still avoid treating those elements as the unknown
quantities of the problem. The unknowns should be the corrections to previously known values,
\textls{approximated as closely as possible}, and so on.

. . . . . . . . . . . .
